% ================================================================
%  chapters/parts/ch07-a.tex
%  第 7 章  原位红外光谱方法（分片 a）
%  源：S1（上册）第 7 章 §7.1，书内 p.413—447 / PDF p.432—466
%  覆盖：§7.1 红外光谱的基本原理和获取原位红外光谱的方法
%        （透射／漫反射／发射／ATR／合频 SFG／时间分辨／近场／原位与 operando）
%  不覆盖：§7.2 及以后（探针分子、金属与氧化物催化剂、反应机理）→ 分片 b/c
% ================================================================

\chapterstart{原位红外光谱方法}%
{本章导读：本章回答——在催化剂真正工作的条件下，能不能、以及怎样从红外光谱读出表面
吸附物种的结构，并据此判断反应机理。第一节建立工具层：红外吸收的条件与定量边界、七种
测量模式各看什么、以及「原位」与「operando」的区别。其后各节转入谱图解释，从探针分子
的红外光谱讲到金属与氧化物催化剂上的反应中间物。分片 a 只覆盖第一节。}
\label{ch:07}

\subsection{原位红外光谱测量什么}

\begin{probout}
红外光谱测量的是\textbf{分子振动跃迁引起的吸收}。落到催化研究上，它给出的是
\textbf{催化剂表面吸附物种的结构信息}：某个物种以什么方式键合（端基或桥式羰基、线性或
桥式吸附、解离或非解离吸附）、键在什么位上，以及这些物种在升温、通气、切换气氛的过程中
如何生成与消失。

\textbf{「原位」（in situ）的含义是：样品在红外池内就地完成焙烧、抽真空脱气、流动氧化
还原、吸附与反应等处理，并在处理与反应条件下就地采谱}，而不是把样品取出后另做表征。
这是红外方法相对动力学方法的根本优势：过去靠对反应物与产物的动力学观察去反推表面中间物，
由于缺少吸附物种结构的直接依据，结论存在相当大的不确定性
\src{S1 §7，书内 p.416 / PDF p.435}。

它\textbf{不}给出：吸附物种的绝对覆盖度（见下文定量边界）、活性位的几何排布、体相结构，
也不能在载体不透明的低波数区给出透射谱。
\end{probout}

红外光谱进入催化研究很早：1930 年 De Boer 研究有机分子在碱金属卤化物上的吸附，
Eischens 等在 1954 年测得 \ce{CO} 在 \ce{Pt}、\ce{Ni} 上吸附的红外光谱后引起广泛兴趣，
此后探针分子的红外光谱成为表征表面「活性位」的常规手段
\src{S1 §7，书内 p.416 / PDF p.435}。

这一方法有四处已知的边界，读后面任何一张原位谱图时都应记住
\src{S1 §7，书内 p.416—417 / PDF p.435—436}：（1）最通用的透射法在研究负载催化剂时，
\textbf{大部分载体在低于约 \qty{1000}{cm^{-1}} 处就不透明了}，这一波数以下的吸附物种很难
测到；（2）金属粒子暴露的表面各不相同（平台、边、角、阶梯、相间界面），同一个吸附物种
可以给出宽化甚至多组的谱带，解释困难；（3）反应中间物的表面浓度低、寿命短，普通傅里叶
变换红外（FTIR）的灵敏度与跟踪速度（毫秒级）往往不够；（4）只有红外活性的物质才有信号，
与红外互补的是拉曼方法。

\subsection{红外吸收的基本原理}

\begin{principle}
\textbf{吸收条件}：只有引起\textbf{偶极矩变化}的振动才能吸收红外光，即
$\left(\partial\mu/\partial Q\right)_{0}\neq0$。因此 \ce{N2}、\ce{O2} 这类同核双原子分子
没有红外吸收（可作载气），而 \ce{H2O}、\ce{CO2} 有很强的红外吸收（是干扰源，见后文）。

\textbf{波数}：$\tilde{\nu}=1/\lambda=\nu/c$，单位 \unit{cm^{-1}}，正比于振动频率，由化学键
强度与折合质量决定。催化研究的主战场是 $4000$--$\qty{400}{cm^{-1}}$ 的中红外区。

\textbf{能量平衡}：固体与电磁波相互作用时，能量按吸收、反射（散射）、透射三部分分配

\begin{equation}\label{eq:art}
A+R+T=1
\end{equation}

式中 $A$、$R$、$T$ 分别为吸收、反射（散射）、透射对总能量的贡献。\textbf{式~\eqref{eq:art}
就是测量模式的选型依据}：样品吸收适当强而散射弱时用透射法；样品散射或反射能量大时用
漫反射法；样品吸收很强（近乎不透明）时用发射法
\src{S1 §7.1，书内 p.417 / PDF p.436}。

\textbf{Beer--Lambert 定律}：吸光度与吸收物种的浓度、光程成正比

\begin{equation}\label{eq:beer}
A=\lg\frac{I_{0}}{I}=\varepsilon\,c\,l
\end{equation}

式中 $\varepsilon$ 为摩尔吸光系数，$c$ 为浓度，$l$ 为光程。
\end{principle}

式~\eqref{eq:beer} 在\textbf{均匀溶液}中是可靠的定量工具，但在\textbf{原位}条件下只有
限地成立，三条原因必须写清：

\begin{itemize}
\item \textbf{消光系数不确定}。吸附态物种的 $\varepsilon$ 随取向、覆盖度、与金属的电子
相互作用而变，通常无法独立测定；同一个物种在载体上与在单晶上的 $\varepsilon$ 不可互借。
\item \textbf{散射损失分不开}。式~\eqref{eq:art} 中的 $R$ 与 $A$ 在粉末样品里无法独立
分离，测到的「吸光度」中混有随波数变化的散射背景，尤其在近红外区散射损失很大。
\item \textbf{光程不确定}。自支撑片的有效厚度不可能精确知道（压片厚度不均），漫反射与
ATR 更是没有一条几何意义上的光程。
\end{itemize}

\begin{keycond}
\textbf{在原位研究中，红外谱带强度只能用于\textbf{同一样品、同一制样、同一条件}下的相对
比较}（随时间的增减、随温度的强弱），不能直接换算成覆盖度或浓度，除非专门标定（如用已知
覆盖度的探针分子、同位素标记或与化学吸附量关联）。峰高仅在 $\varepsilon$ 与光程不变时
与浓度线性；谱带饱和时线性关系首先失效。任何报出的「吸附物种浓度」都必须同时给出标定
方式，否则是不可复核的数值。
\end{keycond}

\subsubsection{傅里叶变换红外光谱：原位研究的默认平台}

现代原位红外几乎都在 FTIR 上完成。其核心是迈克尔逊（Michelson）干涉仪：红外光经准直后在
与光束成 $45^{\circ}$ 的分束器上分成两束，分别射向动镜与定镜，返回后再经分束器合成
\textbf{干涉光}；动镜移动改变光程差 $\delta$，检测器记录的总光强随 $\delta$ 周期变化，
即\textbf{干涉图} $I(\delta)$ \src{S1 §7.1.1，书内 p.422—425 / PDF p.441—444}：

\begin{equation}\label{eq:interferogram}
I(\delta)=\frac{1}{2}\int_{0}^{\infty}B(\tilde{\nu})\left[1+\cos\left(2\pi\tilde{\nu}\delta\right)\right]\mathrm{d}\tilde{\nu}
\end{equation}

式~\eqref{eq:interferogram} 中 $B(\tilde{\nu})$ 是经分束器效率与检测器响应修正后的光源光谱
分布；常数项不携带光谱信息，对干涉图作傅里叶变换即得光谱。采样零点由 He--Ne 激光干涉信号
触发，动镜每移动 \qty{316.45}{nm} 采集一个数据点，光程差间隔因而恒定
\src{S1 §7.1.1，书内 p.423—425 / PDF p.442—444}。

\tocheck{未对照原页面图像：式~\eqref{eq:interferogram} 的系数（$1/2$）、常数项与积分上下限；
OCR 在此残缺（书内 p.423—426 / PDF p.442—445）。He--Ne 波长 OCR 作 632.9 nm（常见值
632.8 nm），此处从 OCR。}

相对色散型（棱镜、光栅）仪器，FTIR 的优势正是原位与瞬态研究的前提
\src{S1 §7.1.1，书内 p.426 / PDF p.445}：分辨率高（\qty{0.1}{cm^{-1}} 不难，商品仪器可达
\qty{0.01}{cm^{-1}} 以下）；\textbf{波数准确度高}（以激光为基准，容易准确到
\qty{0.01}{cm^{-1}} 且不需经常校谱——这正是能分辨几 \unit{cm^{-1}} 吸附位移的原因）；扫描快
（全谱 $1$~s 内完成，最快毫秒级跟踪）；范围宽（$11000$--$\qty{10}{cm^{-1}}$）且信噪比高
（可多次扫描平均）。代价是三者的取舍：\textbf{提高时间分辨率必然降低信噪比}，反之亦然。

\subsection{七种测量模式：按样品与信息深度选型}

红外测量模式的差别，本质上是\textbf{光与样品以什么方式作用}：光穿过体相（透射）、在粉末
表面散射后返回（漫反射）、由样品自身热辐射发出（发射）、在晶体内部全反射而只与表面层耦合
（ATR），以及在此之上的三种特殊实现（合频、时间分辨、近场）。表~\ref{tab:modes} 是本章的
选型骨架：先看第二列的\textbf{光路特点与样品要求}，再看第四列的\textbf{能得到／得不到什么}。

\begin{table}[tbp]
\centering
\caption{七种红外测量模式的选型对照。模式选择由三件事共同决定：\textbf{样品形态}（能否做成
薄片、是否含水、能否压紧）、\textbf{散射强弱}与\textbf{所需信息深度}（体相／表面层／表面
单层）。表中温度为附件与商品池的一般范围，不同厂家差别较大。}
\label{tab:modes}
{\srcsize
\begin{tabularx}{\textwidth}{@{}l >{\raggedright\arraybackslash}X >{\raggedright\arraybackslash}X >{\raggedright\arraybackslash}X@{}}
\toprule
模式 & 光路特点与样品要求 & 适用样品与典型池 & 能得到／不能得到 \\
\midrule
透射吸收 & 入射光穿过样品，测透射光的衰减；样品须对红外透明且足够薄（自支撑片、金属蒸膜、
气溶胶膜），粒径应小于波长以抑制散射 & 可压片的负载催化剂、氧化物、分子筛粉末；真空或
流动气氛的高温池（可到 $800\,^{\circ}\mathrm{C}$） & 得到体相平均的吸附物种与载体羟基信息，
波数范围宽、定量最直接；得不到多数载体 \qty{1000}{cm^{-1}} 以下的谱图，也测不了液体与
致密块体。制样最难、散射损失最大 \\
漫反射 DRIFT & 光斜入射到松散粉末，经颗粒多次散射后返回表面，用离轴椭圆镜收集，避开法线
方向的镜面反射 & 松散粉末、散射很强的样品；原位漫反射池（室温至 $400$--$500\,^{\circ}\mathrm{C}$，
真空或加压），可在线上接质谱、色谱 & 得到表面层物种，制样最简单，可测不透明载体；得不到
可直接互比的定量结果——须经 Kubelka--Munk 变换，且基线随粒度与装填状态漂移 \\
发射 & 不用外光源，直接测样品自身的热辐射，须与同温度的黑体或灰体本底相比 & 高温下工作的
样品（发射池可到 $700\,^{\circ}\mathrm{C}$）、镀金或不锈钢衬底上的薄膜 & 得到强吸收、无法
透射或反射的样品在高温下的谱；得不到高信噪比与可靠的绝对强度，通常只报相对发射度 \\
衰减全反射 ATR & 光在高折射率晶体内部全反射，靠表面倏逝波与样品耦合；要求晶体折射率大于
样品且入射角大于临界角，样品须与晶体紧密接触 & 液体、糊状物、聚合物、致密片；液相原位池、
电化学池、加压接触夹具 & 得到紧贴晶体表面的亚微米层信息，样品形态几乎不受限；得不到体相
信息，穿透深度随波长变化使谱带相对强度失真，粉末必须压紧 \\
红外合频 SFG & 可调红外光与固定可见光同时照射表面，检测和频信号
$\omega_{\mathrm{SF}}=\omega_{\mathrm{IR}}+\omega_{\mathrm{vis}}$，只有非中心对称的表面与
界面才有信号 & 平整表面（单晶、薄膜、气液与液固界面）；$250$--$1000$~K，超高真空至几个
大气压 & 得到不受体相与气相干扰的表面／界面专属信号，可在反应条件下捕捉中间物；得不到
粗糙与多孔样品的谱图，且红外频率范围有限（约 $1100$--$\qty{4000}{cm^{-1}}$） \\
时间分辨 TR-IR & 快扫描（毫秒级）、步进扫描（ns--\textmu s）、泵浦--探测（fs--ps）；
须配快速响应检测器与高速数据采集 & 光催化、瞬态中间物、表面反应动力学；要求所研究的动力学
过程可重复成百上千次 & 得到物种随时间的演化与寿命，是区分中间物与旁观者的关键手段；得不到
同时兼顾高分辨率与高信噪比——提高时间分辨率必然牺牲信噪比 \\
近场 Nano-FTIR & 金属化探针针尖在样品表面形成纳米级增强场，散射光按针尖机械振动频率解调，
只取近场分量 & 表面纳米结构、单个粒子与局域物种分布；须原子力显微镜与 FTIR 联用 & 得到
约 \qty{20}{nm} 空间分辨率的光谱与成像；得不到快速、大面积的测量——扫描慢，针尖与样品
距离控制苛刻，定量很弱 \\
\bottomrule
\end{tabularx}

\vspace{2pt}
{\srcsize 据 S1（上册）§7.1.1—§7.1.7，书内 p.417—441 / PDF p.436—460 归纳；各模式参数为
文献与商品附件的一般范围。}}
\end{table}

\begin{figure}[tbp]
\centering
\begin{tikzpicture}[x=1mm,y=1mm,>=stealth,
  beam/.style={thinline,->},
  weakbeam/.style={thinline,->,draw=black!55},
  scat/.style={thinline,->,draw=black!55,dashed},
  samp/.style={draw=black!70,fill=black!25},
  xtal/.style={draw=black!70,fill=black!10}]
% ---------------- (a) 透射 ----------------
\draw[beam] (2,20) -- (13,20);
\node[anchor=south] at (7.5,20.6) {入射光 $I_{0}$};
\filldraw[samp] (13,11) rectangle (17,29);
\node[anchor=north] at (15,10.4) {样品片};
\draw[weakbeam] (17,20) -- (31,20);
\node[anchor=south] at (24,20.6) {透射光 $I$};
\draw[scat] (17,24.5) -- (26,30);
\draw[scat] (17,15.5) -- (26,10);
\node[anchor=west] at (26.5,30) {散射损失};
\node[anchor=north] at (17,5) {（a）透射：光穿过体相};
% ---------------- (b) 漫反射 ----------------
\filldraw[samp,pattern=north east lines,pattern color=black!50] (62,10) rectangle (102,16);
\node[fill=white,inner sep=1pt] at (82,13) {多次散射与吸收};
\node[anchor=north] at (81,8.6) {松散粉末（不压片）};
\draw[draw=black!60,thinline] (60,31) arc[start angle=90,end angle=270,x radius=3.2mm,y radius=7mm];
\draw[draw=black!60,thinline] (104,31) arc[start angle=90,end angle=-90,x radius=3.2mm,y radius=7mm];
\node[anchor=south] at (60,31.4) {收集镜};
\node[anchor=south] at (104,31.4) {收集镜};
\draw[beam] (46,27) -- (57.1,27);
\draw[beam] (57.1,27) -- (70,16);
\node[anchor=south east] at (46,27.6) {入射光};
\draw[scat] (70,16) -- (78,30);
\draw[scat] (70,16) -- (90,33);
\draw[scat] (70,16) -- (106.9,27);
\draw[beam] (106.9,27) -- (116,27);
\node[anchor=south] at (114,27.6) {到检测器};
\node[anchor=north] at (81,5) {（b）漫反射：光在粉末中散射后返回表面};
% ---------------- (c) ATR ----------------
\filldraw[xtal] (122,18) -- (150,18) -- (158,10) -- (130,10) -- cycle;
\filldraw[samp] (122,21) -- (150,21) -- (150,18) -- (122,18) -- cycle;
\node[anchor=west] at (151.5,21) {样品};
\draw[dotted,thinline,->] (130,18.2) -- (130,21);
\draw[dotted,thinline,->] (146,18.2) -- (146,21);
\node[anchor=west] at (130.8,19.6) {$d_{\mathrm{p}}$};
\draw[beam] (114,6) -- (126,14);
\draw[mainline] (126,14) -- (130,18) -- (138,10) -- (146,18) -- (154,10) -- (156,12);
\draw[beam] (156,12) -- (166,22);
\node[anchor=east] at (113.6,6.6) {入射光};
\node[anchor=west] at (166.4,22) {出射光};
\node[anchor=north] at (140,9.4) {全反射晶体（$n_{1}>n_{2}$）};
\node[anchor=north] at (142,5) {（c）ATR：光在晶体中多次全反射};
\end{tikzpicture}
\caption{三种常用模式中光与样品的相互作用方式对比。读者应看三处：\textbf{光在样品中走的
路径}（穿过体相／在颗粒间多次散射后返回／在晶体内部多次全反射）、\textbf{作用体积}
（体相、表层粉末、表面亚微米层）、以及由此决定的\textbf{样品形态与含水量要求}——透射要求
样品能做成自支撑薄片且散射弱，含水或易碎的粉末难以制样；漫反射直接用松散粉末，不破坏样品
形态，但粉末堆积状态会直接改变基线；ATR 只探测紧贴晶体的表面层，样品形态几乎不受限，但
要求接触紧密、晶体折射率大于样品。图中光路为示意，未按真实比例绘制。\schemamark
\figsrc{据 S1（上册）图 7-10、图 7-14、图 7-16，书内 p.428、p.432、p.433 / PDF p.447、
p.451、p.452 改绘。}}
\label{fig:paths}
\end{figure}

\begin{keycond}
选型的三句口诀（由式~\eqref{eq:art} 直接推出）：\textbf{吸收适中、散射弱 $\rightarrow$ 透射}；
\textbf{散射或反射能量大 $\rightarrow$ 漫反射}；\textbf{吸收很强、近乎不透明 $\rightarrow$
发射}。再叠加两条信息深度要求：只要表面层信息、样品又无法制样时选 ATR；要的是表面单层、
且不允许体相与气相干扰时选 SFG。没有一种池或附件能满足所有研究要求，这也是原位池至今仍多
为自行设计加工的原因 \src{S1 §7.1.1，书内 p.421 / PDF p.440}。
\end{keycond}

\subsubsection{透射红外吸收光谱：默认选项与它的代价}

透射法技术相对简单、适用面广；代价集中在\textbf{制样与散射}，其余条件集中列在下框中。

\textbf{（1）制样与散射}。负载催化剂折射率高，散射是获得好谱图的主要障碍：散射损失取决于样品与
周围介质的折射率差、入射光波长与颗粒大小，因此\textbf{颗粒应小于所用红外波长}以降低散射；
近红外区散射损失很大，远红外区虽散射小但体相吸收大，故一般在 \qty{1000}{cm^{-1}} 以下不
推荐透射测量 \src{S1 §7.1.1，书内 p.418 / PDF p.437}。常用的 KBr 压片与石蜡糊法因稀释剂
吸收与压力效应，\textbf{不能用于原位研究}。原位研究用\textbf{自支撑片}：不加稀释剂直接把
粉末压成薄片，具体条件集中列在下框
\src{S1 §7.1.1，书内 p.418—419 / PDF p.437—438}。

\begin{measure}
\textbf{透射原位的常规测量条件}（据 S1 §7.1.1，书内 p.418—422 / PDF p.437—441）：
（1）\textbf{片}：直径 $15$--$25$~mm 以充满光束，压力 $70$--$110$~atm（约 $7$--$11$~MPa），
加料均匀、必要时垫云母片；
（2）\textbf{厚度判据}：\qty{4000}{cm^{-1}} 处透射率落在 $10\%$--$30\%$——太厚则吸附带饱和
且可利用能量下降，太薄则吸附物种信号太弱；
（3）\textbf{池}：能在池内完成焙烧、流动氧化还原、抽高真空（脱气）、吸附与反应；可随时
移出或移入光路而不受上述处理影响；记录谱图时不受气相组分影响；池本底干扰尽量小；
（4）\textbf{窗口}：\ce{NaCl}、\ce{KBr} 范围宽但潮解；\ce{CaF2}、\ce{BaF2} 耐水却忌铁盐
溶液、对热与振动敏感；\ce{ZnSe} 近年较常用；高温池窗口与法兰需水冷保护；
（5）\textbf{背景}：在同温度、同气氛下采集，并用热电偶直接测样品（或样品架）温度。
\end{measure}

\begin{pitfall}
\textbf{池内样品温度不等于设定温度。}带水冷与外加热的简单池，外壁与中心样品之间存在温差；
高温石英池在温度高时池壁与样品的实际温差最大可达 $100\,^{\circ}\mathrm{C}$
\src{S1 §7.1.1，书内 p.419 / PDF p.438}。因此\textbf{必须用热电偶直接测样品（或样品架）
温度并在文中说明测温方式}；用炉温代替样品温度，会把温度效应误读成物种变化。
\tocheck{未对照原页面图像：简单池「$200\,^{\circ}\mathrm{C}$ 时温差小于 $2\,^{\circ}\mathrm{C}$」
一项，OCR 作「小于 27 ℃」（书内 p.419 / PDF p.438），数值疑为 2 ℃，待核。}
\end{pitfall}

\subsubsection{漫反射红外光谱（DRIFT）：松散粉末与 Kubelka--Munk}

漫反射把「散射」从损失变成信号来源：光斜入射到松散粉末，在颗粒间多次散射后一部分返回表面，
由离轴椭圆镜收集（图~\ref{fig:paths}b）。它可以直接测量松散粉末，避免压片造成的扩散影响，
很适用于散射很强的样品；通常用 DRIFT 或 DRIFTS 表示漫反射傅里叶变换红外光谱。其吸收光谱
用 Kubelka--Munk 函数描述 \src{S1 §7.1.2，书内 p.427 / PDF p.446}：

\begin{equation}\label{eq:km}
\frac{K}{S}=\frac{\left(1-R_{\infty}\right)^{2}}{2R_{\infty}}
\end{equation}

式中 $K$ 为吸收系数（频率的函数），$S$ 为散射系数，$R_{\infty}$ 为无限厚样品的反射比（一般
厚度在几毫米即可满足该条件）。

\begin{keycond}
式~\eqref{eq:km} 把反射比换算成与浓度近似成正比的量，但成立前提是\textbf{散射系数 $S$ 在
测量过程中恒定}——也就是粒度、装填密度、压实程度都不变。这决定了 DRIFT 的使用纪律：
同一系列实验必须保持\textbf{同一批样品、同一粒度、同一装填方式}，背景必须在\textbf{同一
温度与同一气氛}下采集；否则基线漂移会被误当作吸附物种的增减。
\tocheck{未对照原页面图像：式~\eqref{eq:km} 的符号约定（$R_{\infty}$ 为反射比还是反射率、
是否取负对数）与 $K$、$S$ 的定义，OCR 残缺（书内 p.427 / PDF p.446）。}
\end{keycond}

为了做原位研究，已发展出可控气氛与压力的原位漫反射池，可在室温至 $400$--$500\,^{\circ}\mathrm{C}$、
真空或加压条件下工作，并可配气路与产物分析
\src{S1 §7.1.2，书内 p.428 / PDF p.447}。相对透射，DRIFT 的代价是定量更弱、谱带形状受
Kubelka--Munk 变换影响，气相干扰也更难扣除（光路在池内是开放的散射几何）。

\subsubsection{衰减全反射红外光谱（ATR）：只取表面层}

ATR 以光在两种介质界面发生全内反射为基础。当光从光密介质（折射率为 $n_{1}$ 的反射元件）
射向光疏介质（折射率为 $n_{2}$ 的样品），且入射角 $\theta$ 大于临界角 $\theta_{\mathrm{c}}$
时发生全反射：

\begin{equation}\label{eq:crit}
\sin\theta_{\mathrm{c}}=\frac{n_{2}}{n_{1}}\qquad\left(n_{1}>n_{2}\right)
\end{equation}

全反射并非光完全不进入样品：光会在样品表面下一定深度才返回，若该层在特定频率有吸收，反射
光中对应频率就出现吸收带，这就是 ATR 光谱。穿透深度 $d_{\mathrm{p}}$ 定义为辐射波电场强度
衰减至表面处 $1/e$ 的距离：

\begin{equation}\label{eq:dpen}
d_{\mathrm{p}}=\frac{\lambda_{0}}{2\pi\sqrt{n_{1}^{2}\sin^{2}\theta-n_{2}^{2}}}
\end{equation}

式中 $\lambda_{0}$ 为\textbf{真空波长}（约 $2.5$--$25$\,µm），$\theta$ 为入射角。
\tocheck{原书式（书内 p.432 / PDF p.451）写作 $d_{\mathrm{p}}=\lambda_{1}/[2\pi n_{1}
\sqrt{\sin^{2}\theta-(n_{2}/n_{1})^{2}}]$，同时把 $\lambda_{1}$ 定义为「介质内波长」——
若分子已用 $\lambda_{1}=\lambda_{0}/n_{1}$，分母不应再出现 $n_{1}$，原书此处符号定义冲突。
本书按倏逝场振幅 $1/e$ 定义改写为上式的真空波长形式。}
\textbf{数值自检}：$\lambda_{0}=\qty{10}{\micro m}$、$n_{1}=2.4$、$n_{2}=1.5$、
$\theta=45^\circ$ 时 $\sqrt{n_{1}^{2}\sin^{2}\theta-n_{2}^{2}}=\sqrt{0.63}=0.7937$，
得 $d_{\mathrm{p}}=\qty{2.005}{\micro m}$。若照原书写法把 $\lambda_{1}=\lambda_{0}/n_{1}$
再除一次 $n_{1}$，会得到 $\qty{0.8355}{\micro m}$——\textbf{系统性低估穿透深度，误差恰为一个
因子 $n_{1}$}。
由式~\eqref{eq:dpen} 可知：\textbf{穿透深度与真空波长成正比}，波长越长探测越深；入射角
接近临界角时穿透深度急速增大，远大于临界角时变化平缓
\src{S1 §7.1.4，书内 p.432 / PDF p.451}。这带来一个常被忽略的后果：
\textbf{同一张 ATR 谱中不同波数的探测深度不同}，谱带的相对强度因此失真，跨波数比较峰高
必须先作深度校正。

为增强吸收峰强度，现代 ATR 附件多用多重全反射：光在梯形（或平行四边形）晶体中多次全反射，
次数可达 $20$--$50$ 次，总穿透深度大大增加。反射次数由晶体几何决定
\src{S1 §7.1.4，书内 p.433 / PDF p.452}：

\begin{equation}\label{eq:nref}
N=\frac{l}{d}\cot\theta
\end{equation}

式中 $l$ 为全反射晶体的长度，$d$ 为两个反射面间的距离，$\theta$ 为入射角。$l$ 与 $d$ 固定，
减小入射角可增加反射次数，是提高 ATR 信号强度的有效办法。

\tocheck{未对照原页面图像：式~\eqref{eq:dpen} 已按倏逝场振幅 $1/e$ 定义改写为真空波长形式
（见该式下方的说明与数值自检）；式~\eqref{eq:nref} 的 $l$、$d$ 定义与是否需要厚度修正
仍待核对（书内 p.432—433 / PDF p.451—452）。}

晶体材料的选择取决于折射率、化学稳定性、光谱范围与机械强度。由于大多数有机物折射率在
$1.5$ 以下，按 $n_{1}>n_{2}$ 的要求，ATR 晶体折射率须大于 $1.5$，常用 KRS-5、\ce{ZnSe}、
\ce{Ge}、\ce{Si} 等；晶体做成菱形，样品放在两个较大的侧面上，几何尺寸受全反射次数与光源
光斑大小约束 \src{S1 §7.1.4，书内 p.433 / PDF p.452}。其中 \ce{ZnSe} 光谱范围宽、化学
稳定性好，\ce{Ge} 极硬但透光范围窄，KRS-5 质软、有毒且受强酸与烷烃类物质影响。

\begin{pitfall}
\textbf{ATR 的成败几乎全在「接触」。}倏逝波只与紧贴晶体的那一层样品耦合，接触不良时信号
骤降、谱带畸变，且不同波数处接触状态不同会造成假峰。粉末样品的原位气固 ATR 因此受限：
必须加压夹紧，而加压本身又会改变床层孔隙与传质——这一条与下文「压片压力」的误判是同一
类问题。ATR 给出的是表面层信息，\textbf{不能代表体相组成}。
\end{pitfall}

\subsubsection{其余四种模式各自解决什么问题}

\textbf{红外发射光谱}不用外光源，直接测样品自身的热辐射，适用于吸收很强、既透不过也反射
不回来的样品，并天然适合高温条件。依据是 Kirchhoff 定律：热平衡下发射率等于吸收率
（$a=e$），配合 $t+r+a=1$（反射率极小时化为 $e=1-t$），发射光谱可换算为透射光谱。原位
发射池可在池内完成真空处理、氧化还原与反应，样品温度可到 $700\,^{\circ}\mathrm{C}$，衬底
常用镀金不锈钢片以避免高温氧化与反应。因绝对发射率难以测定，实践中以同温度下样品支撑片
（近似灰体）的本底谱 $L_{\mathrm{r}}(T)$ 作参考，报出\textbf{相对发射度}
$RE=L_{\mathrm{s}}(T)/L_{\mathrm{r}}(T)$，可同时扣除支撑片发射与仪器光路的吸收和发射
\src{S1 §7.1.3，书内 p.429—431 / PDF p.448—450}。
\tocheck{未对照原页面图像：相对发射度的定义式 $RE=L_{\mathrm{s}}(T)/L_{\mathrm{r}}(T)$ 与
$e=L_{\mathrm{s}}/L_{\mathrm{bb}}$ 的符号，以及「发射度（emittance）」与「发射率
（emissivity）」两个术语的中文用法（书内 p.431 / PDF p.450）。}

\textbf{红外合频（SFG）}是表面与界面专属的技术。非线性光学过程要求介质非中心对称，一般
物质体相是中心对称的，只有表面与界面处对称性被破坏才产生二次非线性现象；可调红外光与
固定可见光同时照射表面，取和频信号即得表面振动光谱。它的选律是\textbf{振动模式必须同时
具有红外与拉曼活性}，信号不载体相与气相贡献，可在 $250$--$1000$~K、超高真空至几个大气压
范围内工作，从而把超高真空模型研究与真实反应条件连接起来。局限同样明确：\textbf{只能用于
平整表面}，粗糙与多孔表面不能研究；高功率激光可能损伤样品；红外频率范围有限（约
$1100$--$\qty{4000}{cm^{-1}}$）\src{S1 §7.1.5，书内 p.434—436 / PDF p.453—455}。

\textbf{时间分辨红外（TR-IR）}按时间尺度选技术：分钟至小时级用普通稳态采谱；秒至毫秒级用
快速扫描 FTIR（动镜加速，最快几十毫秒完成一次扫描），信噪比的损失可用「同一动力学过程
重复多次、按时间序列累加平均」补回；纳秒至微秒级用步进扫描（动镜停在固定位置采集扰动响应，
再逐点移动）配快速 MCT 检测器，要求过程高度可重复；飞秒至皮秒级用泵浦--探测方案，以光学
延迟平台控制探测光与激发光的时间差。步进扫描保持了 FTIR 的高分辨率与波数准确度，但灵敏度
低于光栅扫描，且难以直接抽取某一波数的动力学曲线
\src{S1 §7.1.6，书内 p.436—440 / PDF p.455—459}。

\textbf{近场光学红外（Nano-FTIR）}突破衍射极限。传统红外成像的分辨率受衍射限制，中红外
波段只有微米量级；把扫描近场光学显微镜与 FTIR 结合，探测金属化探针针尖尖端形成的、尺度与
针尖曲率半径相当的纳米增强场与样品的相互作用，可实现纳米级分辨的光谱采集，目前最高空间
分辨率约 \qty{20}{nm}。探针以频率 $\Omega$ 振动使近场信号被周期调制，对基频或高次谐波解调
即可得到纯近场信号 \src{S1 §7.1.7，书内 p.440—441 / PDF p.459—460}。它的限制是扫描慢、
只测针尖附近的极小区域、定量很弱。

\subsection{原位与 operando：差别必须写清}

这是本章标题的落点，也是最容易被混用的一对词。

\begin{principle}
\textbf{原位（in situ）}：\textbf{只在反应条件下测谱}。样品处于反应温度、压力与气氛中，
光谱给出表面物种「有什么、何时出现、何时消失」。但多数原位池不同时测量反应速率与产物分布，
因此\textbf{无法把表面物种与催化性能直接关联}——当原位谱在真空或非反应条件下测量时，这一
缺陷尤其突出 \src{S1 §7.1.8，书内 p.441 / PDF p.460}。

\textbf{operando}：\textbf{在工作状态下同时进行光谱表征与催化性能测量}，即红外（或其他
谱学）与在线产物分析（质谱、色谱）同时进行，从而得到「反应结构」与「反应中间物」的直接
关系，建立动态构效关系。催化剂表面随环境条件（温度、压力、气相或液相组成）动态变化，
\textbf{只有在反应条件下、并同时知道反应速率，才能谈构效关系}——这就是 operando 存在的
理由 \src{S1 §7.1.8，书内 p.441 / PDF p.460}。
\end{principle}

一句话概括：\textbf{operando $=$ in situ $+$ 同时测催化性能}。原位说「表面上有什么」，
operando 才说「这些物种与反应速率有什么关系」。

\textbf{稳态同位素瞬态动力学分析（SSITKA）}是重要的 operando 红外技术。它在等温等压下快速
用同位素瞬态物种替换其中一种反应物（例如以 \ce{^{13}CO} 替换 \ce{^{12}CO}），同时切换惰性
气体以确定反应器中滞留气体的组成；由于等温等压，催化剂表面在切换过程中没有变化。反应物与
产物用在线色谱定量，瞬态响应物种用质谱记录，通过分析原物种与同位素物种随时间的变化规律，
可得到分子水平的机理信息与中间体的动力学量（浓度、反应活性、速率常数），从而把真正的中间物
与旁观物种区分开 \src{S1 §7.1.8，书内 p.442 / PDF p.461}。

\begin{keycond}
operando 池必须同时满足三组互相牵制的要求，因此设计必然是折中：（1）\textbf{光谱可测}——
窗口透光范围、光程与样品量足以给出可分辨的谱带；（2）\textbf{反应器行为接近真实}——传质
传热、空速、压力、产物滞留体积尽量接近实用反应器；（3）\textbf{产物可定量}——死体积小、
响应快，才能把毫秒级的表面变化与产物信号对齐。报告 operando 结果时应给出池型、空速、样品
量与产物分析的响应时间；缺失这些，表面物种与速率的「同步」就无从判断。
需同时注意两条常见的鸿沟：压力鸿沟（超高真空模型研究 vs 常压／高压真实反应）与材料鸿沟
（单晶与模型薄膜 vs 实际负载催化剂），SFG 一类技术正是为跨越前者而发展起来的。
\end{keycond}

\subsection{从原始谱图到结论}

\begin{dataflow}
（1）\textbf{同条件背景}：在同一温度、同一气氛（最好同一池、同一光程）下采集背景；高温时
还需扣除样品热发射的贡献，否则水、\ce{CO2} 与热辐射都会被读成吸附带。
（2）\textbf{差谱与基线}：以吸附前（或惰性气氛下）的谱作参考做差谱，检查差谱中是否出现
负峰或位置漂移的伪峰——出现即说明背景与样品状态不一致，应重测而不是解释。
（3）\textbf{峰位归属}：以基团特征频率与已知参考谱比对；吸附引起的频率位移（相对气相或
物理吸附态）本身就携带键合方式的信息。
（4）\textbf{变量实验}：改变温度、覆盖度、分压或同位素，看各谱带是同步变化还是有先有后；
\textbf{只有先后次序才能给出动力学信息}。
（5）\textbf{与产物分析关联}：把谱带强度的时间曲线与产物生成速率放在同一时间轴上比较
（operando），同步变化的物种才是候选中间物。
（6）\textbf{定量}：只有在同一系列、同一制样与同一条件下比较峰面积；要给出覆盖度必须说明
标定方式（式~\eqref{eq:beer} 的边界见前文关键条件）。
\end{dataflow}

\begin{obs}
（教学设定）某原位透射实验中，切换反应气后 $2100$--$\qty{2050}{cm^{-1}}$ 区间的谱带在
$30$~s 内明显增强；而池出口质谱给出的产物信号约在 $90$~s 后才达到新的稳态。
\end{obs}

\begin{infer}
谱带增强\textbf{先于}产物响应，说明该物种的生成不依赖产物积累，与「产物二次吸附」这一类
解释不符；但这\textbf{不足以}证明它是催化循环内的中间物——它也可能是与反应无关的旁观物种，
或仅仅是池内气相分子在冷点上的物理吸附。要区分，需要同位素切换（看产物中是否出现同位素
标记）或瞬态实验（看它的消耗速率是否与产物生成速率同量级）。
\end{infer}

\subsection{误判与边界}

\begin{pitfall}
\textbf{（1）原位池的传质与传热滞后。}池内滞留体积大、气体置换需要时间，气相组成的变化
与样品实际感受到的变化之间存在时间差；同时池内存在温度梯度（前文所述的池壁与样品温差）。
后果是：表面物种的响应时间被高估，先后次序可能被颠倒。对策是用质谱或色谱同步监测出口
组成以估计滞后时间，并在同一篇工作中固定气速与样品量。

\textbf{（2）自支撑片的厚度与压片压力改变样品本身。}压力过低片子易碎、厚度不均；压力过高
会把颗粒压实，改变孔隙结构与表面可及性——对分子筛与高比表面氧化物尤其明显，此时测到的
已不是真实催化剂表面的行为。同一系列实验必须固定压力与厚度，并记录片重与厚度。

\textbf{（3）水与 \ce{CO2} 的干扰。}它们是大气与管路中最常见的杂质，气相谱带与表面羟基、
碳酸盐、甲酸盐物种的谱带重叠，极易误判。\supp{补充：\ce{CO2} 的反对称伸缩振动在约
\qty{2349}{cm^{-1}}，水的弯曲振动在约 \qty{1600}{cm^{-1}}、伸缩振动在
$3400$--$\qty{3600}{cm^{-1}}$ 附近，都是原位谱图中最常出现的干扰位置。}对策是充分吹扫或
抽真空、并在同条件下取背景；但要注意\textbf{不能一律扣除}——反应中真实生成的表面碳酸盐、
甲酸盐物种恰恰就在这些位置，把它们当干扰扣掉会删掉真正的结论。

\textbf{（4）「看到谱带变化」不等于「看到催化循环内的中间物」。}谱带出现或消失只说明表面
浓度变了，并没有检验该物种是否与催化循环的速率相关。判据至少要有其一：它与产物生成的
速率同量级并可建立定量关系；同位素切换或瞬态实验给出预期响应；预吸附该物种会改变（或不
改变）转化频率。缺少这些证据，只能称为「反应条件下的表面物种」，不能称为中间物。

\textbf{（5）升温带来的峰位移与热发射。}温度升高时键长与吸附强度变化会使峰位移动，高温下
样品的热发射还会叠加入透射谱形成假峰。跨温度比较峰位时，应先用同一温度下的背景扣除，并
避免把几十 \unit{cm^{-1}} 的整体位移解释为「生成了新物种」。
\end{pitfall}

\begin{crosscheck}
红外只给振动信息，必须与其他方法互证才能支撑结论。（1）\textbf{与拉曼互补}：红外活性源于
偶极矩变化，拉曼活性源于极化率变化，中心对称分子遵循互斥规则；载体在
\qty{1000}{cm^{-1}} 以下不透光的问题可用拉曼绕开，两者对同一物种的归属一致时才可确信。
（2）\textbf{与程序升温脱附、化学吸附、质谱联用}：给出吸附量与脱附温度，与红外的物种归属
相互印证；活性位数目的定量属于化学吸附方法的范围（\chapref{ch:05}）。
（3）\textbf{与 X 射线衍射等体相方法联用}：区分表面物种与体相变化（\chapref{ch:04}）。
（4）\textbf{与在线色谱／质谱联用}构成 operando，把表面物种与反应速率关联；催化剂比表面
与孔结构的变化则回到物理吸附（\chapref{ch:01}）。
若红外与拉曼、或红外与脱附实验给出矛盾结论，应先检查标定、基线与样品状态，而不是择一采信。
\end{crosscheck}

\srcnote{本分片来源：S1（上册）§7.1（含 §7.1.1—§7.1.8），书内 p.413—447 / PDF p.432—466。
其中式~\eqref{eq:art}、式~\eqref{eq:km}、式~\eqref{eq:crit}--\eqref{eq:nref} 的变量定义按
正文 OCR 重建，系数与符号约定存疑处已在文中以「待核对」标出；图~\ref{fig:paths} 为据原书
图 7-10、7-14、7-16 改绘的示意光路，非原图描摹。§7.2 及以后的探针分子与反应机理内容
见分片 b/c。}
